% ABSTRACT -- first full draft; numbers marked PENDING are filled once E2/E4-E6/E8 land.
% Order: problem -> gap -> what we do -> strongest results -> meaning (5 sentences before any detail).
\begin{abstract}
Protein interaction networks and knowledge graphs are now a standard ingredient of models that predict
transcriptional responses to drugs and genetic perturbations, and published ablations credit the network
with substantial gains.
Those ablations compare against random graphs that destroy the degree distribution, so they cannot tell
whether a model uses \emph{which} genes interact or only \emph{how many} interactions each gene has.
We introduce a control ladder for network priors---degree-preserving rewiring, node-label permutation,
edge-count-matched random graphs and an empty graph---together with a model-free information test that
measures what a network can carry about responses before any model is trained.
Across eight interaction networks and three perturbation datasets (LINCS L1000 at two phases and Tahoe-100M),
the dominant network signal is target degree: drugs whose targets are hubs elicit stronger responses
(Spearman $\rho = 0.12$--$0.34$), an association that survives correction for publication bias, while
wiring beyond degree is small in undirected protein networks and appears clearly only in directed
signalling and transcription-factor networks (+0.03 to +0.07).
Auditing published and new models with the ladder, the real network is indistinguishable from its
degree-preserving rewiring in DeepCE (gene vectors replaceable by a one-hot identity), in graph neural
networks over STRING, and \pending{in GSNN, GEARS and PDGrapher}, whereas the drug-side information that
does transfer to unseen compounds is target identity, morphology and measured responses.
The ladder is a cheap, model-agnostic standard for claiming that a network prior helps, and it points to
directed regulatory structure as the network information worth modelling.
\end{abstract}

\keywords{perturbation response \and protein interaction networks \and graph neural networks \and null models \and LINCS L1000 \and Tahoe-100M}
